%% Appendix D: Near-ceiling stress-tests and supporting corpora
%% Demoted from Chapter 7 main spine (breakout iteration).
%% Most labels preserved so existing \ref{} targets keep working.

\chapter{Near-Ceiling Stress-Tests and Supporting Corpora}
\label{app:stress}

This appendix holds material that supports---but must not lead---the
conjunctive-release argument.
Main-text claims rest on E6, E2, equal-$K$, C4, and the ACL interception case
(Chapter~\ref{ch:results}).
Numbers below are fixed-oracle unless an archive/pre-fix caption says
otherwise; do not mix corpora in one claim sentence.

%=============================================================================
\section{Quality--Overhead Archive (Pareto / Latency)}
% label kept on Chapter pointer: sec:results:rq5
%=============================================================================

Archive Pareto/latency figures below use pre-fix exports where captioned;
authoritative cost is the fixed-oracle one-line card in Chapter~\ref{ch:results}.

\begin{figure}[t]
  \centering
  \widefig{figures/mode_distribution_latency.pdf}
  \caption{\textbf{Archive figure (pre-fix E1 export):} historical latency
    mass under the buggy others-witness oracle / hard gate.
    \textbf{Authoritative cost story} is fixed-oracle E1
    (\texttt{run\_e1\_m\_win\_v2}): M 10.5\,s vs.\ B2 9.6\,s at comparable
    mean llm\_calls (${\approx}1.16$ vs.\ ${\approx}1.18$)---do not cite this
    figure as the current premium.
    \textbf{How to read:} log-scale wall-clock latency per task on the
    historical export only.}
  \label{fig:latency-distribution}
\end{figure}

The latency distribution (Figure~\ref{fig:latency-distribution}) shows repair
modes paying a clear premium: M and A1 spread toward higher wall-clock times,
while B1/B2 remain near one-call latencies.

\begin{figure}[t]
  \centering
  \widefig{figures/pareto_quality_vs_latency.pdf}
  \caption{\textbf{Claim (latency axis from pre-fix Pareto export):}
    B2 is the cheaper mean-Conf.\ / latency operating point; M buys
    prevention / conjunctive Accept at ${\approx}3.3\times$ latency
    (46.2\,s vs.\ B2 13.9\,s on that export).
    Primary fixed-oracle Conf.\ is M~100.0\% / B2~98.3\%
    (Table~\ref{tab:main-results})---do not read the figure's pre-fix
    Conf.\ annotations (B2~87.7\%, M~86.3\%) as current primary quality.
    \textbf{How to read:} right/up is better quality, left is cheaper; the
    dashed polyline is the Pareto frontier and the annotation marks mode~M.}
  \label{fig:pareto}
\end{figure}

The insight in Figure~\ref{fig:pareto} is a trade-off, not a free lunch.
On the pre-fix latency export, M pays $3.3\times$ vs.\ B2
(46.2\,s vs.\ 13.9\,s), mostly extra LLM calls (2.08 vs.\ 1.24)---acceptable
for parallel CI, not for interactive IDE use.
Historical A2 (39.9\,s, 80.3\% Conf.\ on pre-fix E1) traded prevention for
modest latency savings vs.\ that run's M; pre-fix B2 (13.9\,s, 87.7\%) was
the cheaper mean-Conf.\ point on that export.
Under the fixed oracle, primary quality is M~100.0\% vs.\ B2~98.3\% at
comparable llm\_calls (Table~\ref{tab:main-results}); the $3.3\times$ latency
premium remains the cost framing for deploying full~M when Accept/PDR justify
it.
Mann-Whitney confirms the latency gap is not a sampling artefact
($p{<}0.001$).

\paragraph{Answer to RQ5.}
Best prevention (PDR/FAR) at ${\approx}3.3\times$ B2 latency (pre-fix cost
export); default to B2 for mean Conf.\ / latency and deploy M when
Accept/PDR-FAR justify the cost.
Fixed-oracle E1 aggregate ranking (Section~\ref{sec:results:e1-detail}) is the
stress-test companion to this Pareto reading---not a claim that M dominates
B2 on every corpus.

%=============================================================================
\section{E1 Aggregate Stress-Test}
\label{sec:results:e1-detail}
%=============================================================================

\textbf{Claim (C2--C4 jointly; RQ5 stress-test).}
The full loop should raise mean Conf.\ over one-shot and
test-feedback baselines; under the fixed others-witness oracle with
strengthened~M ($K{=}5$, best-effort selection, advisory critical-only gate),
Strict tracks Conf.\ closely on BENCH-120.
Aggregate E1 ranking is reported here as an overlap stress-test under RQ5,
not as a universal-effectiveness headline across unfiltered corpora.

\paragraph{Equal-$K$ hygiene (cross-ref).}
Primary equal-$K$ Conf.\ ranking is Table~\ref{tab:equal-k} in
Section~\ref{sec:results:rq4} (\texttt{M\_eq} ${+}2.5$\,pp vs.\ B2).
Lead mechanism evidence remains E6 (within~M, $K{=}3$, ${+}7.7$\,pp with CI
excluding zero).
Strengthened E1 M-win ($K{=}5$ bundle) below is a near-ceiling stress-test.

\subsection{Main Results}

The headline question for this E1 stress-test is whether specification-guided
repair improves \emph{conformance} under a fair first-match witness oracle.
Table~\ref{tab:main-results} answers with B1/B2/M on 120 tasks each from
\texttt{run\_e1\_m\_win\_v2}; ablation and extended-baseline rows from the
pre-fix canonical run are retained only as historical context
(footnote).

\begin{table}[H]
  \caption{Main comparison results ($n = 120$ tasks per mode).
    \textbf{B1/B2/M} from fixed-oracle \texttt{run\_e1\_m\_win\_v2}
    (strengthened~M: best-effort $\arg\max$ Conf., E14-informed
    \texttt{execution\_trace\_matched}, advisory critical-only pattern gate,
    $K{=}5$, \texttt{formal\_max\_cases}${}{=}32$).
    Strict = strict success / acceptance rate
    ($\mathrm{Accept}{=}1$); Conf.\ = mean formal conformance
    ($\overline{\mathrm{Conf}}$); Kill = mutation kill rate (checker precision on
    injected mutants, not generator quality---see \S\ref{sec:results:e1-detail}); Fail.\ = mean
    strict evaluation failures per task.
    \emph{Note:} Primary B1/B2/M use the fixed first-match others-witness
    oracle (\texttt{run\_e1\_m\_win\_v2}).
    Fixed-oracle component ablations (A1--A3) are Table~\ref{tab:ablation-fixed}
    (\texttt{run\_e1\_ablation\_fixed\_v1}).
    B3--B5 and daggered A1--A3 rows below are historical
    \texttt{run\_e1\_canonical\_v1} under the buggy others-witness oracle and
    are \emph{not} primary claims.
    Report M vs.\ B2 as win rate and means (2/120 wins, 0 losses, 118 ties;
    ${+}1.7$\,pp); no Holm claim is asserted for this run.
    Equal-$K$ Conf.\ ranking is Table~\ref{tab:equal-k} (\texttt{M\_eq} vs.\ B2).}
  \label{tab:main-results}
  \centering
  \thesistable{%
    \footnotesize
    \renewcommand{\arraystretch}{1.12}
    \begin{tabular*}{\linewidth}{@{\extracolsep{\fill}}l r r r r r@{}}
      \toprule
      Mode & Strict & Conf. & Kill & Lat.\ (ms) & Fail. \\
      \midrule
      B0$^\dagger$   & 5.0\%           & 89.2\%        & 48.6\%       & 165           & 0.00 \\
      B1   & 84.2\%          & 84.2\%        & 52.4\%       & 17{,}193      & 0.00 \\
      B2   & 98.3\%          & 98.3\%        & 49.1\%       & 9{,}576       & 0.00 \\
      B3$^\dagger$   & 5.0\%           & 81.9\%        & 51.0\%       & 13{,}210      & 2.33 \\
      B4$^\dagger$   & 5.0\%           & 87.1\%        & 49.3\%       & 10{,}758      & 2.49 \\
      B5$^\dagger$   & 5.0\%           & 81.1\%        & 50.6\%       & 6{,}212       & 2.33 \\
      M    & \textbf{100.0\%} & \textbf{100.0\%} & 48.6\%    & 10{,}457      & 0.00 \\
      A1$^\dagger$   & 5.0\%           & 85.5\%        & 54.0\%       & 4{,}791       & 0.00 \\
      A2$^\dagger$   & 5.0\%           & 80.3\%        & 49.1\%       & 6{,}405       & 0.00 \\
      A3$^\dagger$   & 3.3\%           & 82.4\%        & 53.3\%       & 4{,}333       & 0.00 \\
      \bottomrule
      \multicolumn{6}{l}{\footnotesize $^\dagger$Historical pre-fix
        (\texttt{run\_e1\_canonical\_v1}; not primary).} \\
    \end{tabular*}%
  }
\end{table}

Table~\ref{tab:main-results} states the aggregate means under the fixed
oracle: M reaches 100\% Conf./Strict, ahead of B2 (98.3\%) and B1 (84.2\%).
Primary component ranking is the fixed-oracle ablation
(Table~\ref{tab:ablation-fixed}: A3 $-$25.8\,pp, A1 $-$17.5\,pp, A2 $0.0$\,pp);
daggered historical A1--A3 rows are archive only.

\begin{figure}[t]
  \centering
  \widefig{figures/mode_distribution_strict_conformance.pdf}
  \caption{\textbf{Claim (pre-fix export):} M concentrates mass in the upper
    conformance band; B0 does not.
    \textbf{How to read:} each violin/box is per-task $\mathit{Conf}(c,t)$ over
    120 tasks; compare the right-hand mass of M with B0's low mode.
    Primary fixed-oracle means are in Table~\ref{tab:main-results}.}
  \label{fig:conf-distribution}
\end{figure}

Figure~\ref{fig:conf-distribution} separates difficulty calibration from
LLM quality on the pre-fix export: B0's bimodal spread marks template
failure on hard tasks, while B1--M cluster in the 60--100\% band.

\begin{figure}[t]
  \centering
  \widefig{figures/formal_conformance_cdf.pdf}
  \caption{\textbf{Claim (pre-fix export):} B2 and A2 dominate the upper CDF
    tail on the historical corpus.
    \textbf{How to read:} curves farther right mean higher conformance for more
    tasks; B0's steep early rise marks template failure on hard tasks.
    Primary fixed-oracle means are in Table~\ref{tab:main-results}.}
  \label{fig:conf-cdf}
\end{figure}

The CDF in Figure~\ref{fig:conf-cdf} is from the pre-fix export: B2 and A2
occupy the upper tail on that corpus; do not treat it as the fixed-oracle
ranking (M~100.0\% $>$ B2~98.3\%).

\begin{figure}[t]
  \centering
  \widefig{figures/summary_metrics_by_mode.pdf}
  \caption{\textbf{Claim (pre-fix export):} highest mean Conf.\ coincides with
    the most LLM calls on that corpus.
    \textbf{How to read:} left axis = Strict / Conf.; right axis = mean LLM
    calls; M's 2.08 calls show the repair loop is exercised, not idle.
    Fixed-oracle llm\_calls are ${\approx}1.16$ (M) vs.\ ${\approx}1.18$ (B2).}
  \label{fig:summary-metrics}
\end{figure}

Figure~\ref{fig:summary-metrics} links quality to cost on the pre-fix
export; under the fixed oracle, M and B2 have comparable mean llm\_calls
near one call (Table~\ref{tab:main-results}).

\begin{figure}[t]
  \centering
  \widefig{figures/llm_calls_distribution.pdf}
  \caption{\textbf{Claim:} budget $K=3$ is actively used, not decorative
    (pre-fix call-count export).
    \textbf{How to read:} M and A1 mass at 2--3 calls; B1 stays at one.}
  \label{fig:llm-calls}
\end{figure}

The call-count histogram (Figure~\ref{fig:llm-calls}) confirms that budget
$K{=}3$ is exercised on the pre-fix export: M and A1 concentrate at two or
three calls, whereas B1 remains at a single generation attempt.

\begin{figure}[t]
  \centering
  \widefig{figures/radar_metrics.pdf}
  \caption{\textbf{Claim (pre-fix export):} M trades latency for conformance;
    B2 is the cheaper moderate alternative on that corpus.
    \textbf{How to read:} axes are normalised Conf., Strict, kill rate, and
    inverse latency; larger area is not uniformly better.
    Primary fixed-oracle Conf.\ ranking is Table~\ref{tab:main-results}.}
  \label{fig:radar-metrics}
\end{figure}

Figure~\ref{fig:radar-metrics} summarises the multi-metric trade-off on the
pre-fix export; fixed-oracle primary quality remains M~100.0\% vs.\ B2~98.3\%.

Table~\ref{tab:main-results} is the aggregate on fixed-oracle E1: M at
100.0\% Conf.\ / Strict leads B2 (98.3\%) and B1 (84.2\%).
Against B2, M wins 2/120 paired tasks (${+}1.7$\,pp; losses~0, ties~118);
mean llm\_calls are comparable (${\approx}1.16$ vs.\ ${\approx}1.18$).
We report win rate and means only---no Holm significance is claimed for this
run.
That is the honest primary-quality reading: structured specification feedback
matters in controlled E6 (+7.7\,pp vs.\ unstructured test-only under full~M;
E14: length-matched \texttt{semantic\_ir} 75.1\% does not match
\texttt{execution\_trace\_matched} 85.4\%) and E11 deployment slices;
fixed-oracle E10/E12 place M and B2 near ceiling without M dominating every
seed.
Historical B3--B5 / A1--A3 rows under the buggy others-witness oracle
(${\approx}5\%$ Strict) are not primary claims.

\subsection{Reading the Distributions}

Three readings settle the rest of the E1 stress-test.

\textbf{Strict now tracks Conf.}
Under the fixed oracle, B1/B2/M Strict matches Conf.\ (84.2 / 98.3 /
100.0\%): the historical uniform ${\approx}5\%$ Strict on
\texttt{run\_e1\_canonical\_v1} was an others-witness measurement artefact,
not the current primary result.

\textbf{The M--B2 margin is real but thin.}
M leads B2 by ${+}1.7$\,pp with only 2/120 decisive wins (118 ties);
mean llm\_calls stay near one call (${\approx}1.16$ vs.\ ${\approx}1.18$).

\textbf{Mutation kill is about the checker, not the generator.}
Kill rates of ${\approx}49$--$52\%$ across modes still measure checker
\emph{precision} against injected syntactic faults, not generator quality.

\paragraph{E1 stress-test note (RQ5).}
Fixed-oracle E1: M 100.0\% $>$ B2 98.3\% $>$ B1 84.2\% (win rate 2/120 vs.\ B2;
no Holm claim).
E6 confirms typed feedback mechanism (+7.7\,pp vs.\ unstructured test-only
under full~M; E14 scopes against length-matched traces).
Prevention (PDR/FAR) remains the complementary discriminator for conjunctive
release; fixed-oracle E10/E12 (Sections~\ref{sec:results:e10},~\ref{sec:results:e12})
bound seed stability and unfiltered ranking.

%=============================================================================
\section{Boundary Density (E4)}
\label{sec:results:e4}
%=============================================================================

Test-feedback looks strong until witness sampling is held constant.
If denser random tests already covered first-match regions, SMT witnesses
would be optional; Figure~\ref{fig:boundary} demonstrates they are not.

\begin{figure}[t]
  \centering
  \widefig{figures/boundary_coverage_by_density.pdf}
  \caption{\textbf{Claim:} random sampling collapses in the Dense tier; SMT
    does not.
    \textbf{How to read:} coverage vs.\ case budget by density; compare SMT
    vs.\ random within the Dense row at budget~$\ge 8$.}
  \label{fig:boundary}
\end{figure}

At budget 16, random sampling covers far less of the Dense tier than
SMT-guided generation; the gap shrinks in the Sparse tier.
That is why test-feedback baselines look adequate on simple specifications
and degrade on hard ones---the same story as the overlap-tertile split in
Section~\ref{sec:results:rq2}, now measured at the sampling layer alone.

%=============================================================================
\section{Generalisation Study (E8)}
\label{ch:generalization}
%=============================================================================

External validity asks whether SpecIR adapters carry the prevention loop
beyond SOFL syntax, or whether the gains are notation-specific.
Table~\ref{tab:generalisation} and Figure~\ref{fig:generalisation} answer
that question.

\begin{figure}[t]
  \centering
  \widefig{figures/generalisation_bar.pdf}
  \caption{\textbf{Claim:} structured feedback helps across notations, but
    adapter transfer is partial on the expanded $n{=}30$ corpus.
    \textbf{How to read:} bars are mean Conf.\ by notation and mode; compare
    M vs.\ B1 within each group (M leads on both adapters), then M vs.\ B2.}
  \label{fig:generalisation}
\end{figure}

On SOFL/FSF (fixed-oracle E1), M leads at 100.0\% versus B2 98.3\% and B1 84.2\%.
On the expanded adapter corpora ($n{=}30$ each, \texttt{run\_e8b\_expanded\_v1}),
M beats B1 on Mini-StateMachine (41.6\% vs.\ 24.7\%) and Mini-Z (23.3\% vs.\
5.8\%); M is within 0.8~pp of B2 on Mini-StateMachine and leads B2 by 5.0~pp
on Mini-Z.
Absolute conformance on adapter notations remains low (23--42\%), reflecting
notation-specific generation difficulty rather than SOFL-level overlap structure;
HumanEval/MBPP rows in Table~\ref{tab:generalisation} show stronger transfer on
real-derived guards.
Chapter~\ref{ch:discussion} interprets these boundaries.

\begin{table}[H]
  \caption{\textbf{Key observation:} M leads B1 on all notations; adapter
    absolute conformance is low on the expanded $n{=}30$ corpus (E8b).
    HumanEval-FSF and MBPP-FSF rows report full B1/B2/M from
    \texttt{run\_e8c\_full\_v1}.}
  \label{tab:generalisation}
  \centering
  \thesistable{%
    \small
    \begin{tabular}{lccc}
      \toprule
      Notation & B1 Conf.\% & B2 Conf.\% & M Conf.\% \\
      \midrule
      SOFL/FSF (120 tasks) & 84.2 & 98.3 & \textbf{100.0} \\
      Mini-StateMachine (30 tasks) & 24.7 & 40.8 & \textbf{41.6} \\
      Mini-Z (30 tasks) & 5.8 & 18.3 & \textbf{23.3} \\
      HumanEval-FSF (20 tasks) & 89.7 & \textbf{98.8} & 87.1 \\
      MBPP-FSF (20 tasks) & 90.0 & 95.0 & 95.0 \\
      \bottomrule
    \end{tabular}%
  }
\end{table}

%=============================================================================
\subsection{External Validity: Real-Derived Benchmark (E8c)}
\label{subsec:results:real-derived}
%=============================================================================

Table~\ref{tab:benchmark-by-source} reports conformance across benchmark subsets.
Mode~M reaches 95.0\% on MBPP-FSF (matching B2) and 87.1\% on HumanEval-FSF
(where B2 leads at 98.8\%).
MBPP confirms partial transfer to real-world guard structure; HumanEval shows
that test-feedback baselines can dominate when branches are sparse.

\begin{inlinetable}{Conformance by benchmark source (mode M, $K=3$).}
  \label{tab:benchmark-by-source}
  \footnotesize
  \begin{tabular}{lrrr}
    \toprule
    Subset & $n$ & M Conf (\%) & M Strict (\%) \\
    \midrule
    Synthetic hard (E1)          & 120 & 100.0 & 100.0 \\
    HumanEval-FSF (E8c)          & 20  & 87.1 & 70.0 \\
    MBPP-FSF (E8c)               & 20  & 95.0 & 85.0 \\
    Mini-Z (E8b expanded)        & 30  & 23.3 & --- \\
    Mini-StateMachine (E8b)      & 30  & 41.6 & --- \\
    \bottomrule
  \end{tabular}
\end{inlinetable}

The higher strict-success rate on real-derived tasks (70--85\%) versus the
historical pre-fix synthetic Strict (${\approx}5\%$) reflected the buggy
others-witness oracle as much as simpler HumanEval/MBPP guard structure;
under the fixed oracle, synthetic hard E1 Strict for~M is 100.0\%
(Table~\ref{tab:main-results}).
Full B1/B2/M baselines on HumanEval-FSF and MBPP-FSF are reported in
Table~\ref{tab:generalisation} from \texttt{run\_e8c\_full\_v2}.

%=============================================================================
\section{Random Benchmark (E10)}
\label{sec:results:e10}
%=============================================================================

E10 tests whether aggregate M beats B2 on an unfiltered random sample
(Section~\ref{sec:exp:e10}).
Table~\ref{tab:e10-random} reports fixed-oracle results from
\texttt{run\_e10\_m\_win\_v1} ($n{=}100$ per mode).
On this \emph{unfiltered} sample, M reaches 100.0\% mean Conf.\ vs.\ B2
99.0\% (${+}1.0$\,pp; wins 1/100, losses~0, ties~99) and B1 81.0\%---consistent
with E1's near-ceiling M$\ge$B2 reading, not a large unfiltered gap.

\begin{table}[H]
  \caption{Random benchmark results (E10, fixed others-witness oracle).
    Uniformly sampled tasks without Z3 overlap filter; modes B1, B2, M;
    \texttt{run\_e10\_m\_win\_v1}.
    M leads B2 by ${+}1.0$\,pp (1/100 wins).}
  \label{tab:e10-random}
  \centering
  \thesistable{%
    \small
    \begin{tabular}{lccc}
      \toprule
      Mode & Conf.\ (\%) & Win vs.\ B2 & Notes \\
      \midrule
      B1 & 81.0 & --- & --- \\
      B2 & 99.0 & ref. & --- \\
      M  & \textbf{100.0} & 1/100 & ${+}1.0$\,pp; 0 losses, 99 ties \\
      \bottomrule
    \end{tabular}%
  }
\end{table}

The decision criterion remains qualitative: fixed-oracle E10 places M slightly
ahead of B2 near ceiling; Figure~\ref{fig:e10-overlap} tertiles should be read
together with E3 (Table~\ref{tab:e3-tiers})---neither shows monotone M
dominance by overlap alone.
The honest claim is overlap-\emph{conditional} deployment (C4), with E6 as the
mechanism anchor.
(Pre-fix \texttt{run\_e10\_random\_v2} had favoured B2 89.1\% vs.\ M 83.1\%
under the buggy others-witness oracle and is not primary.)

\begin{figure}[t]
  \centering
  \widefig{figures/performance_vs_overlap.pdf}
  \caption{E10: mean conformance by overlap-rate tertile on the unfiltered random
    benchmark (figure from pre-fix export; read aggregate means from
    Table~\ref{tab:e10-random}).}
  \label{fig:e10-overlap}
\end{figure}

%=============================================================================
\section{External SOFL Benchmark (E11)}
\label{sec:results:e11}
%=============================================================================

E11 evaluates curated, non-synthetic FSF tasks with citable external provenance
(Section~\ref{sec:exp:e11}).
Table~\ref{tab:e11-external} reports repeat~0 results from
\texttt{run\_e11\_external\_v1} ($n{=}22$; zero ID overlap with E1).

\begin{table}[H]
  \caption{External SOFL benchmark (E11).  Hand-curated from textbook and IET
    pattern sources; not generated by \texttt{HardTaskGenerator}.}
  \label{tab:e11-external}
  \centering
  \thesistable{%
    \small
    \begin{tabular}{lcccc}
      \toprule
      Mode & $n$ & Strict (\%) & Conf.\ (\%) & Lat.\ (ms) \\
      \midrule
      B1 & 22 & 63.6 & 89.8 & 2{,}018 \\
      B2 & 22 & 63.6 & 89.8 & 4{,}171 \\
      M  & 22 & 59.1 & 81.9 & 5{,}103 \\
      \bottomrule
    \end{tabular}%
  }
\end{table}

On the external corpus under the primary endpoint (\texttt{ecnu-plus}),
B1 and B2 tie at 89.8\% mean conformance; M trails
(81.9\%)---these curated tasks have moderate guard overlap and simpler branch
structure than the E1 hard set, consistent with the deployment boundary in
Chapter~\ref{ch:discussion}.
Replicating E11 on \texttt{ecnu-max} (\texttt{run\_e11\_ecnu\_max\_v1}) yields
a three-way tie at 89.8\% Conf.\ / 63.6\% Strict for B1/B2/M: the stronger
endpoint closes the primary-model M gap without reversing B2 as the default
mean-Conf.\ choice.
Provenance for each task is recorded in \texttt{benchmarks/external\_sofl.json}.
E15 stratification (overlap-density tiers on the external corpus) explains why
M trails B1/B2 on \texttt{ecnu-plus}: curated tasks have moderate overlap and
simpler branch structure than the E1 stress-test
(see \texttt{e11\_overlap\_stratified.csv}).

%=============================================================================
\section{Multi-Seed Stability (E12)}
\label{sec:results:e12}
%=============================================================================

E12 addresses LLM stochasticity on the \textbf{full} 120-task hard benchmark
(Section~\ref{sec:exp:e12}).
Table~\ref{tab:e12-stability} reports fixed-oracle E12-full from
\texttt{run\_e12\_m\_win\_v1} (120 tasks $\times$ seeds $\{0,1,2\}$ $\times$
modes B1, B2, M; 1080 jobs).
Overall means: B1 88.6\%, B2 100.0\%, M 99.7\%.
B2 is first every seed (ties with M at 100\% on seeds~0 and~2); M's sole
failure is seed~1 HardCase069 (Conf.\ 0).
Mean-across-seeds pairing: M wins 0, losses~1, ties~119 vs.\ B2.
Honest framing: both near ceiling; E1/E10 show M$\ge$B2 on aggregate, while
E12 shows B2 slightly more seed-stable (100.0\% vs.\ 99.7\%)---do \emph{not}
claim M dominates all seeds.

\begin{table}[H]
  \caption{E12-full multi-seed stability under fixed others-witness oracle
    (\texttt{run\_e12\_m\_win\_v1}; 120 tasks, seeds $\{0,1,2\}$).
    Per-seed mean formal conformance (\%).}
  \label{tab:e12-stability}
  \centering
  \thesistable{%
    \small
    \begin{tabular}{lcccc}
      \toprule
      Mode & Mean Conf.\ (\%) & Seed 0 & Seed 1 & Seed 2 \\
      \midrule
      B1 & 88.6 & 95.0 & 75.8 & 95.0 \\
      B2 & \textbf{100.0} & 100.0 & 100.0 & 100.0 \\
      M  & 99.7 & 100.0 & 99.2 & 100.0 \\
      \addlinespace
      \multicolumn{5}{l}{\textit{Mean-across-seeds M vs.\ B2: wins 0, losses 1, ties 119}} \\
      \multicolumn{5}{l}{\textit{Ranking: B2 first every seed (ties M @100\% on seeds 0, 2)}} \\
      \bottomrule
    \end{tabular}%
  }
\end{table}

\subsection{Second-Model Replication (E16)}
\label{sec:results:e16}

E16 replicates B1/B2/M on the \textbf{full} 120-task BENCH-120 using
\texttt{ecnu-max} (DeepSeek-family) while E1/E12 use \texttt{ecnu-plus}
(Qwen-family) (\texttt{run\_e16\_canonical\_v1}, $n{=}120$ per mode).
All three modes tie at \textbf{89.2\%} mean strict conformance---endpoint
choice does not reverse the deployment boundary
(Table~\ref{tab:deployment-boundary}), though absolute gaps shrink vs.\ the
primary model.
The controlled E6 protocol on the same endpoint
(Table~\ref{tab:feedback-e6-max}) likewise saturates (${\approx}89.2\%$ for
A/B/C), confirming that typed-IR gains require repair headroom.
The prior 29-task pilot (\texttt{run\_e16\_model\_pilot\_v1}) is superseded.

\begin{table}[H]
  \caption{E16 second-model replication (\texttt{ecnu-max}, $n{=}120$).}
  \label{tab:e16-canonical}
  \centering
  \thesistable{%
    \small
    \begin{tabular}{lcc}
      \toprule
      Mode & Conf.\ (\%) & Lat.\ (ms) \\
      \midrule
      B1 & 89.2 & 3{,}374 \\
      B2 & 89.2 & 4{,}139 \\
      M  & 89.2 & 10{,}072 \\
      \bottomrule
    \end{tabular}%
  }
\end{table}

%=============================================================================
\section{Commercial Cross-Model Replication (n1n)}
\label{sec:results:n1n}
%=============================================================================

Campus endpoints alone leave a residual single-gateway threat.
We replicate the controlled E6 A/B/C sweep and B1/B2/M ranking on commercial
models via the n1n OpenAI-compatible gateway
(\url{https://api.n1n.ai/pricing}): \texttt{gpt-4o},
\texttt{claude-sonnet-4-6}, and \texttt{deepseek-v3.2}, on the same
29-task stratified subset as E12
(\texttt{benchmarks/e12\_stratified\_30.json}).

% Auto-generated by aggregate_n1n_campaigns.py — do not edit by hand.
\begin{table}[t]
  \centering
  \caption{E6 feedback variants across n1n models (A=\texttt{test\_only}, B=\texttt{test\_expected}, C=\texttt{semantic\_ir}). Conf.\ as mean formal conformance (\%); $\Delta$ is C$-$A in percentage points.}
  \label{tab:cross-model-e6}
  \footnotesize
  \begin{tabular}{lcccc}
    \toprule
    Model & A (\%) & B (\%) & C (\%) & $\Delta$(C$-$A) \\
    \midrule
    gpt-4o & 88.9 & 87.9 & 88.9 & 0.0 \\
    claude-4.6 & 87.5 & 81.0 & \textbf{88.5} & 1.0 \\
    deepseek & 75.0 & 69.8 & 71.1 & -3.9 \\
    \bottomrule
  \end{tabular}
\end{table}


\paragraph{Reading Table~\ref{tab:cross-model-e6}.}
On \texttt{gpt-4o}, A and C tie at 88.9\% (saturation, $\Delta{=}0$).
On \texttt{claude-sonnet-4-6}, typed IR yields a modest
\textbf{+1.0\,pp} (87.5$\rightarrow$88.5).
On \texttt{deepseek-v3.2}, C trails A by 3.9\,pp (75.0$\rightarrow$71.1):
longer structured prompts can hurt weaker commercial endpoints.
Together with primary E6 (+7.7\,pp on \texttt{ecnu-plus}) and E6-max
saturation (${+}0.03$\,pp), the mechanism claim is headroom- and
model-dependent, not universal.

% Auto-generated by aggregate_n1n_campaigns.py — do not edit by hand.
\begin{table}[t]
  \centering
  \caption{E16 stratified subset (B1/B2/M) on n1n endpoints.}
  \label{tab:cross-model-e16}
  \footnotesize
  \begin{tabular}{llccc}
    \toprule
    Model & Mode & $n$ & Conf.\ (\%) & Strict (\%) \\
    \midrule
    gpt-4o & B1 & 29 & 100.0 & 0.0 \\
     & B2 & 29 & 100.0 & 0.0 \\
     & M & 29 & 88.9 & 0.0 \\
    \bottomrule
  \end{tabular}
\end{table}


On the same stratified hard subset, \texttt{gpt-4o} ranks
B1${=}$B2 at 100\% Conf.\ with M at 88.9\%---again no case for replacing B2
as the mean-Conf.\ default.

%=============================================================================
\section{Industrial SOFL Corpus}
\label{sec:results:industrial}
%=============================================================================

\subsection{Real-Priority Micro-Benchmark}
\label{sec:results:real-priority}

Table~\ref{tab:real-priority-micro} reports equal-$K$ B1/B2/\texttt{M\_eq}
on the curated real-priority micro-benchmark
(Section~\ref{sec:exp:real-priority}): firewall / role / billing / tax
adaptations plus industrial, GitHub, HKCA09, and published-industrial
slices---not HardSynthetic tasks.
Reading follows C4: if mean Conf.\ ties or saturates, the reason to deploy
full~M remains acceptance safety (Accept / FAR), not a Conf.\ crown; if B1
leaves headroom, the corpus still shows that the pipeline runs on
practitioner-style ordered guards outside the lab generator.

% Auto-generated by summarize_real_priority_micro.py
\begin{table}[t]
  \centering
  \caption{\textbf{Real-priority micro-benchmark}
    (\texttt{real\_priority\_micro\_v1}, $n{=}30$; equal $K{=}3$,
    \texttt{ecnu-plus}).
    Curated ACL / billing / routing / firewall / role-gate / tax-bracket
    tasks from industrial patterns, GitHub \texttt{.asfl} harvest, HKCA09,
    and published-industrial pilots---not HardSynthetic generators.
    Lead reading: external validity for C4 (Accept/FAR deployment), not a
    Conf.\ crown over B2.}
  \label{tab:real-priority-micro}
  \footnotesize
  \renewcommand{\arraystretch}{1.12}
  \begin{tabular}{@{}lcccc@{}}
    \toprule
    Mode & Mean Conf.\ (\%) & Strict (\%) & vs.\ B2 (W/L/T) & Mean llm\_calls \\
    \midrule
    B1 & 95.6 & 93.3 & --- & 1.00 \\
    B2 & 100.0 & 100.0 & --- & 1.07 \\
    \texttt{M\_eq} & 100.0 & 100.0 & 0/0/30 & 1.00 \\
    \bottomrule
  \end{tabular}
\end{table}


\subsection{SMT Domain Scalability}
\label{sec:results:smt-scalability}

Table~\ref{tab:smt-scalability} and Figure~\ref{fig:smt-domain-latency}
answer the ``Z3 only works on $[-5,20]$'' concern for the guards that
matter here.
On \texttt{real\_priority\_micro\_v1}, mean witness time stays
${\approx}30$\,ms from $[-5,20]$ through $[-1000,1000]$; a Z3
\texttt{Real} encoding of the same predicates is even slightly cheaper
(${\approx}18$\,ms).
In short: \emph{widening the LIA box does not explode} for these
first-match linear guards.
Stress probes with more variables or a mild nonlinear product remain
sub-second on this hardware
(\texttt{artifacts/smt\_scalability\_v1/stress\_probe.json}).
When solvers would time out or return \texttt{unknown} on richer theories,
Section~\ref{sec:disc:smt-hybrid} specifies the hybrid fuzzing fallback
rather than claiming universal SMT scalability.

% SMT domain scalability on real_priority_micro_v1 (n=30).
\begin{table}[t]
  \centering
  \caption{\textbf{SMT scalability ablation}
    (witness generation on \texttt{real\_priority\_micro\_v1}, $n{=}30$,
    3 repeats). Widening the LIA box from the historical default
    $[-5,20]$ to $[-1000,1000]$ does \emph{not} explode latency on these
    first-match linear guards; a Z3 \texttt{Real} encoding of the same
    predicates is likewise cheap. Stress probes (more variables /
    nonlinear products) grow but stay sub-second here
    (Figure~\ref{fig:smt-domain-latency};
    \texttt{artifacts/smt\_scalability\_v1/}).}
  \label{tab:smt-scalability}
  \footnotesize
  \renewcommand{\arraystretch}{1.12}
  \begin{tabular}{@{}lcccc@{}}
    \toprule
    Domain / encoding & Mean (ms) & p50 & p95 & Max \\
    \midrule
    LIA $[-5,20]$ (default) & 29.1 & 31.0 & 39.1 & 46.7 \\
    LIA $[-20,50]$ & 28.9 & 33.0 & 37.9 & 38.7 \\
    LIA $[-100,100]$ & 29.6 & 32.8 & 39.6 & 44.6 \\
    LIA $[-1000,1000]$ & 30.5 & 34.0 & 39.9 & 43.7 \\
    Real $[-100,100]$ & 18.6 & 19.2 & 24.8 & 32.0 \\
    \bottomrule
  \end{tabular}
\end{table}


\begin{figure}[t]
  \centering
  \widefig{figures/smt_domain_latency.pdf}
  \caption{\textbf{Claim:} for FSF-style linear first-match guards on the
    real-priority micro-benchmark, widening the integer domain (and probing
    Z3 Reals) does not produce a latency cliff.
    \textbf{How to read:} solid $=$ mean ms/task; dashed $=$ p95.
    Nonlinear / high-dimensional stress is discussed in
    Section~\ref{sec:disc:smt-hybrid}, not hidden.}
  \label{fig:smt-domain-latency}
\end{figure}

Beyond textbook/IET external SOFL (E11), we evaluate a curated
industrial-pattern corpus
(\texttt{benchmarks/industrial\_sofl.json}, $n{=}31$): banking, telecom,
access control, medical alerts, SCADA-like thresholds, insurance, and related
ordered-guard decision tables, with provenance
\texttt{liu2004soflbook} / \texttt{li2023iet\_faultprevention} /
\texttt{industrial\_pattern:<domain>}
(\emph{not} production vendor traces; see
\texttt{benchmarks/INDUSTRIAL\_SOFL\_README.md}).

% Auto-generated by aggregate_n1n_campaigns.py — do not edit by hand.
\begin{table}[t]
  \centering
  \caption{Industrial SOFL corpus (B1/B2/M) on n1n endpoints. Conf.\ = mean formal conformance; Strict = conjunctive accept rate.}
  \label{tab:cross-model-industrial}
  \footnotesize
  \begin{tabular}{llcccc}
    \toprule
    Model & Mode & $n$ & Conf.\ (\%) & Strict (\%) & Lat.\ (ms) \\
    \midrule
    gpt-4o & B1 & 31 & 95.5 & 67.7 & 20670 \\
     & B2 & 31 & 100.0 & 67.7 & 16809 \\
     & M & 31 & 95.5 & 67.7 & 32401 \\
    \addlinespace
    claude-4.6 & B1 & 31 & 95.5 & 67.7 & 19253 \\
     & B2 & 31 & 100.0 & 67.7 & 20479 \\
     & M & 31 & 95.5 & 67.7 & 21565 \\
    \addlinespace
    deepseek & B1 & 31 & 52.7 & 16.1 & 7946 \\
     & B2 & 31 & 80.7 & 41.9 & 10330 \\
     & M & 31 & 72.5 & 32.3 & 14941 \\
    \bottomrule
  \end{tabular}
\end{table}


\paragraph{Reading Table~\ref{tab:cross-model-industrial}.}
Under both \texttt{gpt-4o} and \texttt{claude-sonnet-4-6}, test-feedback B2
reaches \textbf{100\%} mean Conf.\ while B1 and M sit at 95.5\%, with
identical Strict 67.7\%.
On \texttt{deepseek-v3.2} absolute Conf.\ is lower, but the ranking is
unchanged: B2 (80.7\%) $>$ M (72.5\%) $>$ B1 (52.7\%).
This is the strongest external-validity support for the deployment rule:
on practitioner-style ordered-guard specs, prefer B2 for mean conformance;
reserve full~M when conjunctive $\mathrm{Accept}$ or PDR/FAR gates are required.

\paragraph{Published-industrial / desensitized real SOFL pilot.}
Beyond pattern corpora, we evaluate $n{=}28$ FSF tasks
\emph{reconstructed and desensitized} from published industrial SOFL case
studies---railway crossing~\cite{liu1998railwaycrossing}, interlocking
\cite{luo2017railway}, Agile-SOFL ATM constraints~\cite{li2022qrs_companion},
and book applications~\cite{liu2004soflbook,li2023iet_faultprevention}
(\texttt{benchmarks/published\_industrial\_pilot.json};
\texttt{run\_desens\_real\_sofl\_v1}; Z3 validation 0 fails).
Under equal $K{=}3$ on \texttt{ecnu-plus}: B1/B2/\texttt{M\_eq} all reach
\textbf{100.0\%} mean Conf.\ / Strict (paired M\_eq vs.\ B2: 0/0/28)
(Table~\ref{tab:desens-real-sofl}).
This is \textbf{not} a proprietary Casco/Mitsubishi dump; vendor
\texttt{.asfl} files merge when placed under
\texttt{vendor/agile-sofl-toolchain/examples} (see \texttt{vendor/README.md}).
Saturation under a strong endpoint reinforces C4: mean Conf.\ alone does not
justify escalating from B2 on published-industrial ordered-guard tasks.

% Desensitized published-industrial SOFL pilot (Wave-2, n=28).
\begin{table}[t]
  \centering
  \caption{Desensitized published-industrial SOFL pilot
    (\texttt{run\_desens\_real\_sofl\_v1}; $n{=}28$ tasks; equal $K{=}3$;
    \texttt{ecnu-plus}).
    Reconstructed from public SOFL/interlocking/Agile-SOFL sources
    (Z3 validation 0 fails); \emph{not} proprietary vendor dumps.
    All modes saturate---supports C4 default B2 for mean Conf.}
  \label{tab:desens-real-sofl}
  \footnotesize
  \begin{tabular}{lccc}
    \toprule
    Mode & Mean Conf.\ (\%) & Strict (\%) & vs.\ B2 (W/L/T) \\
    \midrule
    B1 & 100.0 & 100.0 & --- \\
    B2 & 100.0 & 100.0 & --- \\
    \texttt{M\_eq} & 100.0 & 100.0 & 0/0/28 \\
    \bottomrule
  \end{tabular}
\end{table}


\paragraph{GitHub live harvest (Wave-3).}
Using authenticated GitHub search we downloaded public formal-spec artefacts
and auto-extracted FSF processes from Agile-SOFL \texttt{.asfl} examples
(\texttt{JaredYe04/agile-sofl-toolchain}; also discovered
\texttt{HKCA09/SOFL-Maintainability-Experiment-Dataset} and
\texttt{ICARUSxyz/FCoTFL}).
After Z3 validation, $n{=}48$ tasks form
\texttt{benchmarks/github\_harvest\_v1.json}
(\texttt{run\_github\_harvest\_v1}).
Under equal $K{=}3$ on \texttt{ecnu-plus}: B1 85.8\%, B2 89.9\%,
\texttt{M\_eq} 89.9\% (paired M\_eq vs.\ B2 1/1/46)
(Table~\ref{tab:github-harvest}).
Unlike the saturated published-industrial pilot, this public harvest leaves
headroom---and still selects B2 (tied with \texttt{M\_eq}) over B1 for mean
Conf., reinforcing C4.

% GitHub live harvest evaluation (Wave-3 refresh).
\begin{table}[t]
  \centering
  \caption{GitHub live harvest corpus
    (\texttt{run\_github\_harvest\_v2}; $n{=}48$; equal $K{=}3$;
    \texttt{ecnu-plus}).
    Auto-extracted from public Agile-SOFL \texttt{.asfl} examples
    (ecommerce, smart-city, hospital, library; Z3-validated).
    Not proprietary vendor dumps. B1 retains headroom; B2 and \texttt{M\_eq}
    tie on mean Conf.\ (0/0/48)---Accept$\neq$mean Conf.; supports C4 default B2.}
  \label{tab:github-harvest}
  \footnotesize
  \begin{tabular}{lcccc}
    \toprule
    Mode & Mean Conf.\ (\%) & Strict (\%) & vs.\ B2 (W/L/T) \\
    \midrule
    B1 & 83.7 & 83.7 & --- \\
    B2 & \textbf{92.0} & 92.0 & --- \\
    \texttt{M\_eq} & 92.0 & 92.0 & 0/0/48 \\
    \bottomrule
  \end{tabular}
\end{table}


\paragraph{HKCA09 SOFL modules $\rightarrow$ FSF (expansion).}
Full SOFL texts from
\texttt{HKCA09/SOFL-Maintainability-Experiment-Dataset}
(ATM, course/hospital registration, vending, stock, transit) are not
auto-encodable as SMT FSF (maps/sequences).
We reconstruct $n{=}35$ ordered-guard integer processes that preserve
decision precedence
(\texttt{benchmarks/hkca09\_sofl\_fsf.json}; RealSpec total $203$).
Equal-$K$ ranking (\texttt{run\_hkca09\_b1b2m\_v1}): B1 74.3\%,
B2/\texttt{M\_eq} 100\% (0/0/35)---B1 headroom with B2 sufficient for
\textbf{C4} (Table~\ref{tab:hkca09-ranking}).
One-shot E6 on a 57-task real/headroom slice saturates at Conf.\ $=1$ for
all feedback variants: aggregate Conf.\ on this slice cannot carry C2.
Hard-seed repair under \texttt{ecnu-plus}
(\texttt{run\_hkca09\_hard\_seed\_e6\_v1}) is \emph{directional} on
\texttt{wrong\_relop} (semantic\_ir 82.0\% vs.\ test\_only 69.3\%;
${+}12.7$\,pp; 9/5/12; CI includes~0) and mixed on structural seeds
(Table~\ref{tab:hkca09-hard-seed}).
We therefore treat HKCA09 as C4 external validity plus a scoped hard-seed
probe---\textbf{not} a second C2 headline.
The publishable C2 advantage remains BENCH-120 E6 (${+}7.7$\,pp, CI
excludes~0) and hard combo-seed support on \texttt{gemini-2.5-flash}
(Table~\ref{tab:gemini-combo-n80}).

% HKCA09 public SOFL FSF reconstructions (Wave-3 expansion).
\begin{table}[t]
  \centering
  \caption{HKCA09 public SOFL modules reconstructed as ordered-guard FSF
    (\texttt{run\_hkca09\_b1b2m\_v1}; $n{=}35$; equal $K{=}3$;
    \texttt{ecnu-plus}).
    Desensitized integer guards preserve auth/capacity/balance precedence
    from the maintainability-experiment materials (ATM, registration,
    vending, stock, transit). Not proprietary dumps.
    B1 has headroom (9/35 fail); B2 and \texttt{M\_eq} both reach 100\%
    Conf.\ --- supports C4 default B2 on this slice.}
  \label{tab:hkca09-ranking}
  \footnotesize
  \begin{tabular}{lcccc}
    \toprule
    Mode & Mean Conf.\ (\%) & Strict (\%) & vs.\ B2 (W/L/T) \\
    \midrule
    B1 & 74.3 & 74.3 & --- \\
    B2 & \textbf{100.0} & 100.0 & --- \\
    \texttt{M\_eq} & 100.0 & 100.0 & 0/0/35 \\
    \bottomrule
  \end{tabular}
\end{table}

% HKCA09 hard-seed feedback ablation (scoped external probe; not C2 headline).
\begin{table}[t]
  \centering
  \caption{Hard-seed repair on HKCA09 FSF under \texttt{ecnu-plus}
    (\texttt{run\_hkca09\_hard\_seed\_e6\_v1}; freeze buggy code $\rightarrow$
    one $T{=}0$ repair).
    One-shot E6 on the same corpus saturates at Conf.\ $=1$.
    \texttt{wrong\_relop} is directional (${+}12.7$\,pp; CI includes~0);
    structural seeds are mixed.
    \textbf{Do not} treat this table as a second C2 headline---primary C2
    remains E6 and Table~\ref{tab:gemini-combo-n80}.
    On \texttt{gemini-2.5-flash}, the same \texttt{wrong\_relop} protocol
    favors test-only (negative control; omitted from the lead).}
  \label{tab:hkca09-hard-seed}
  \footnotesize
  \begin{tabular}{lcccc}
    \toprule
    Seed & test\_only & semantic\_ir & $\Delta$ (pp) & W/L/T \\
    \midrule
    \texttt{wrong\_relop} & 69.3 & \textbf{82.0} & $+12.7$ & 9/5/12 \\
    \texttt{invert\_order} & 96.3 & 93.9 & $-2.3$ & 2/6/22 \\
    \bottomrule
  \end{tabular}
\end{table}


\paragraph{Industrial decision procedure (C4).}
Escalate to~M \emph{iff} the release bar requires $\mathrm{Accept}{=}1$
\emph{or} a prevention target FAR${\le}5\%$ (E2 operating point) \emph{and}
repair headroom remains; otherwise keep B2.
Do \emph{not} escalate from overlap tertiles alone (E3 non-monotone).
On industrial-pattern and desensitized published-industrial pilots, mean-Conf.\
alone already selects B2 under strong endpoints.

%=============================================================================
\section{VerifierLoop-FSF Baseline (E13 / B6)}
\label{sec:results:b6}
%=============================================================================

REV-7 implements mode~B6: SMT witnesses during repair with structured
counterexample prompts but without Semantic Feedback IR or pattern guard
(Section~\ref{sec:exp:methods}).
Table~\ref{tab:b6-stratified} compares B6 with B2 and M on the same 30-task
stratified subset as E12 (\texttt{run\_b6\_stratified\_v1}, repeat~0).
B6 (86.0\% mean conformance) matches B2 (86.2\%) within 0.2~pp while M trails
(83.3\%) on this subset---showing that witness-guided repair without the full IR
recovers test-feedback-level aggregate performance, but does not close the gap to
M on the full hard benchmark (Table~\ref{tab:b6-full}).

\begin{table}[H]
  \caption{VerifierLoop-FSF (B6) vs.\ B2/M on 30-task stratified subset (E13).
    Same tasks and $K{=}3$ as E12; single repeat.}
  \label{tab:b6-stratified}
  \centering
  \thesistable{%
    \small
    \begin{tabular}{lccc}
      \toprule
      Mode & $n$ & Conf.\ (\%) & Lat.\ (ms) \\
      \midrule
      B2 & 30 & 86.2 & 5{,}416 \\
      B6 & 30 & 86.0 & 16{,}245 \\
      M  & 30 & 83.3 & 15{,}099 \\
      \bottomrule
    \end{tabular}%
  }
\end{table}

\begin{inlinetable}{B6 full hard-benchmark run (\texttt{run\_b6\_full\_v2}, $n{=}120$).}
  \label{tab:b6-full}
  \footnotesize
  \begin{tabular}{lrrrr}
    \toprule
    Mode & Strict (\%) & Conf.\ (\%) & Kill (\%) & Lat.\ (ms) \\
    \midrule
    B2 & 5.0 & 88.2 & 61.4 & 7{,}189 \\
    B6 & 5.0 & 81.1 & 51.0 & 21{,}440 \\
    M  & 5.0 & 82.0 & 60.6 & 22{,}507 \\
    \bottomrule
  \end{tabular}
\end{inlinetable}

On the paired full 120-task export (\texttt{run\_b6\_full\_v2},
Table~\ref{tab:b6-full}), B2 reaches 88.2\% Conf., B6 81.1\%, and M 82.0\%.
Thus any SMT counterexample in the loop is \emph{not} enough for an aggregate
Conf.\ win: B6 ${\approx}$~M, and both trail B2.
Typed Semantic Feedback IR's value is shown under the controlled E6 setting
(and in conjunctive $\mathrm{Accept}$ / PDR--FAR), not as ``M beats B6 on mean
Conf.''
B6 latency ($\approx$21\,s) exceeds B2 ($\approx$7\,s) but stays near this run's M
($\approx$23\,s), consistent with formal checking inside the loop without
pattern screening.

\subsection{M\_lite: Witness + Semantic IR (E18)}
\label{sec:results:m-lite}

Mode \texttt{M\_lite} combines B6-style SMT witnesses in the repair loop with
Semantic Feedback IR but \emph{without} the pattern guard
(\texttt{run\_m\_lite\_v1}, $n{=}120$).
Table~\ref{tab:m-lite} answers whether B2$+$typed IR is enough: M\_lite
(79.0\%) trails B6 (81.8\%) and that run's B2 (89.3\%)---naive stacking fails.
The E6 +7.7\,pp effect (vs.\ unstructured test-only under full~M) requires the
full mode~M repair integration (IR$+$hard gate), not witnesses plus IR alone;
practitioners should choose B2 (default mean Conf./latency) or full~M
(Accept/PDR at ${\approx}3.3\times$ latency), not M\_lite as a shortcut.

\begin{table}[H]
  \caption{M\_lite vs.\ B6 vs.\ B2 on 120-task hard benchmark (E18).}
  \label{tab:m-lite}
  \centering
  \thesistable{%
    \small
    \begin{tabular}{lccc}
      \toprule
      Mode & Conf.\ (\%) & Strict (\%) & Lat.\ (ms) \\
      \midrule
      B2      & 89.3 & 5.0 & 9{,}066 \\
      B6      & 81.8 & 4.2 & 36{,}747 \\
      M\_lite & 79.0 & 5.0 & 37{,}219 \\
      M (E1)$^\dagger$  & 86.3 & 5.0 & 46{,}161 \\
      \bottomrule
      \multicolumn{4}{l}{\footnotesize $^\dagger$Historical pre-fix M
        (\texttt{run\_e1\_canonical\_v1}); primary fixed-oracle M is 100.0\%.} \\
    \end{tabular}%
  }
\end{table}

%=============================================================================
\section{Sensitivity Analysis}
\label{sec:results:sensitivity}
%=============================================================================

The sensitivity question is whether the reported operating point is brittle.
Figure~\ref{fig:sensitivity-heatmap} maps mean strict formal conformance against
guard-overlap density and repair budget across modes.

\begin{figure}[t]
  \centering
  \widefig{figures/sensitivity_heatmap.pdf}
  \caption{Sensitivity heatmap: mean strict formal conformance as a function of
    task complexity (guard-overlap tertile) and pipeline configuration.
    Overlap tertiles are non-monotone for M vs.\ B2 on canonical E3; read with
    Table~\ref{tab:e3-tiers}.}
  \label{fig:sensitivity-heatmap}
\end{figure}

\subsection{Sensitivity to Task Complexity}

The 120-task benchmark spans a range of complexity in terms of guard-overlap density
and scenario count.
Partitioning tasks into three tertiles based on guard-overlap rate shows
non-monotone M vs.\ B2 ranking on canonical E3 (Table~\ref{tab:e3-tiers});
higher overlap does not uniformly favour M.
The sensitivity heatmap in Figure~\ref{fig:sensitivity-heatmap} visualises the
same heterogeneity across all ten primary and ablation modes.

\subsection{Sensitivity to Budget $K$}

The budget parameter $K$ controls the maximum number of LLM calls in the repair
loop.
Mode A3 ($K=1$ effectively) achieves 82.4\% conformance on the
\emph{historical} pre-fix corpus; full mode~M ($K=3$) reaches 86.3\% on
pre-fix canonical E1 at seed~0---repair budget alone does not close the
B1/B2 gap without overlap-conditional deployment.
(Primary fixed-oracle E1: M~100.0\% / B2~98.3\%; A1--A3 not re-run.)
The marginal return of the second iteration (roughly 3--4\,pp) is higher than that of
the third (roughly 2--3\,pp), consistent with diminishing returns
(Figure~\ref{fig:convergence}).
Extrapolating, $K=2$ would likely yield conformance in the range 87--88\% on that
historical corpus, comparable to pre-fix B2 but with Z3-guided feedback.
The $K=3$ budget appears to be a reasonable operating point: it captures most of the
attainable gain while keeping total LLM calls below 3.0 per task on average.

\subsection{Sensitivity to Formal Case Count}

The formal checker generates up to 64 conformance cases per task for final strict
scoring, but uses only 16 cases per iteration during the repair loop.
The 16-case internal budget trades coverage for latency; the 64-case strict pass
improves sensitivity to guard-overlap faults on the final reported score.
Increasing the internal case count from 16 to 32 would improve counterexample precision
at the cost of approximately 50\% more Z3 time per iteration---a trade-off identified
for future empirical investigation.

\subsection{Sensitivity to LLM Choice}

All experiments used Qwen2.5-27B-Instruct (\textsf{Qwen-27B}).
The experimental design is not specific to this model; the only model-specific
assumptions are (a)~the model can follow a structured Python code generation prompt,
and (b)~the model can incorporate counterexample inputs into a repair prompt.
A stronger base model would likely improve all LLM-based modes
proportionally, with the B1/M gap potentially widening if that model learns to
avoid ordering errors from the first attempt.
Conversely, a weaker model would benefit more from the repair loop, increasing M's
relative advantage.

%=============================================================================
\section{Statistical Test Summary}
\label{sec:results:stats}
%=============================================================================

\begin{table}[t]
\caption{Statistical test summary for E1 (120-task hard benchmark, repeat~0). Wilcoxon signed-rank for per-task strict success and conformance; Mann--Whitney~U for latency. Holm--Bonferroni correction within each metric family. Cliff's $\delta$: positive favours the first mode in each pair.}
\label{tab:stats}
\centering
\small
\begin{tabular}{lccc}
\toprule
Comparison & $p$ (Holm) & Significant ($\alpha{=}0.05$) & Cliff's $\delta$ / ratio \\
\midrule
\multicolumn{4}{l}{\textit{Primary comparisons}} \\
\addlinespace
\texttt{M} vs \texttt{B1} (strict success) & 1.0000 & No & -0.0167 \\
\texttt{M} vs \texttt{B2} (strict success) & 1.0000 & No & -0.0167 \\
\texttt{M} vs \texttt{B1} (conformance) & 0.0745 & No & +0.0358 \\
\texttt{M} vs \texttt{B2} (conformance) & 0.8113 & No & +0.0002 \\
\texttt{M} vs \texttt{B2} (latency) & <0.001 & Yes & 4.05$\times$ \\
\addlinespace
\multicolumn{4}{l}{\textit{All pairwise conformance (Holm-corrected)}} \\
\addlinespace
\texttt{A1} vs \texttt{B0} (Conf) & <0.001 & Yes & +0.4750 \\
\texttt{A1} vs \texttt{B1} (Conf) & 1.0000 & No & +0.0227 \\
\texttt{A1} vs \texttt{B2} (Conf) & 1.0000 & No & -0.0112 \\
\texttt{A1} vs \texttt{M} (Conf) & 1.0000 & No & -0.0119 \\
\texttt{A2} vs \texttt{A1} (Conf) & 0.9931 & No & +0.0285 \\
\texttt{A2} vs \texttt{B0} (Conf) & <0.001 & Yes & +0.5112 \\
\texttt{A2} vs \texttt{B1} (Conf) & 0.0745 & No & +0.0517 \\
\texttt{A2} vs \texttt{B2} (Conf) & 1.0000 & No & +0.0171 \\
\texttt{A2} vs \texttt{M} (Conf) & 1.0000 & No & +0.0174 \\
\texttt{A3} vs \texttt{A1} (Conf) & 1.0000 & No & -0.0272 \\
\texttt{A3} vs \texttt{A2} (Conf) & 1.0000 & No & -0.0568 \\
\texttt{A3} vs \texttt{B0} (Conf) & <0.001 & Yes & +0.4569 \\
\texttt{A3} vs \texttt{B1} (Conf) & 1.0000 & No & -0.0044 \\
\texttt{A3} vs \texttt{B2} (Conf) & 1.0000 & No & -0.0387 \\
\texttt{A3} vs \texttt{M} (Conf) & 1.0000 & No & -0.0407 \\
\texttt{B1} vs \texttt{B0} (Conf) & <0.001 & Yes & +0.4569 \\
\texttt{B2} vs \texttt{B0} (Conf) & <0.001 & Yes & +0.4871 \\
\texttt{B2} vs \texttt{B1} (Conf) & 1.0000 & No & +0.0341 \\
\texttt{M} vs \texttt{B0} (Conf) & <0.001 & Yes & +0.5067 \\
\texttt{M} vs \texttt{B1} (Conf) & 0.0745 & No & +0.0358 \\
\texttt{M} vs \texttt{B2} (Conf) & 0.8113 & No & +0.0002 \\
\addlinespace
\multicolumn{4}{l}{\textit{All pairwise strict success (Holm-corrected)}} \\
\addlinespace
\texttt{A1} vs \texttt{B0} (strict) & 1.0000 & No & -0.0083 \\
\texttt{A1} vs \texttt{B1} (strict) & 1.0000 & No & +0.0000 \\
\texttt{A1} vs \texttt{B2} (strict) & 1.0000 & No & +0.0000 \\
\texttt{A1} vs \texttt{M} (strict) & 1.0000 & No & +0.0167 \\
\texttt{A2} vs \texttt{A1} (strict) & 1.0000 & No & +0.0000 \\
\texttt{A2} vs \texttt{B0} (strict) & 1.0000 & No & -0.0083 \\
\texttt{A2} vs \texttt{B1} (strict) & 1.0000 & No & +0.0000 \\
\texttt{A2} vs \texttt{B2} (strict) & 1.0000 & No & +0.0000 \\
\texttt{A2} vs \texttt{M} (strict) & 1.0000 & No & +0.0167 \\
\texttt{A3} vs \texttt{A1} (strict) & 1.0000 & No & +0.0000 \\
\texttt{A3} vs \texttt{A2} (strict) & 1.0000 & No & +0.0000 \\
\texttt{A3} vs \texttt{B0} (strict) & 1.0000 & No & -0.0083 \\
\texttt{A3} vs \texttt{B1} (strict) & 1.0000 & No & +0.0000 \\
\texttt{A3} vs \texttt{B2} (strict) & 1.0000 & No & +0.0000 \\
\texttt{A3} vs \texttt{M} (strict) & 1.0000 & No & +0.0167 \\
\texttt{B1} vs \texttt{B0} (strict) & 1.0000 & No & -0.0083 \\
\texttt{B2} vs \texttt{B0} (strict) & 1.0000 & No & -0.0083 \\
\texttt{B2} vs \texttt{B1} (strict) & 1.0000 & No & +0.0000 \\
\texttt{M} vs \texttt{B0} (strict) & 1.0000 & No & -0.0250 \\
\texttt{M} vs \texttt{B1} (strict) & 1.0000 & No & -0.0167 \\
\texttt{M} vs \texttt{B2} (strict) & 1.0000 & No & -0.0167 \\
\bottomrule
\end{tabular}
\end{table}


The decision implication of Table~\ref{tab:stats} is metric choice.
Strict-success comparisons among LLM modes are not significant after
Holm correction ($p{=}1.000$): the conjunctive gate is selective by design.
Latency comparisons confirm the expected cost of M over B2
(Mann--Whitney $p{<}0.001$, $4.05\times$ slower).
Aggregate conformance gaps on pre-fix canonical E1 were not Holm-significant
(M vs.\ B2 $p{=}1.0$); fixed-oracle E1 reports M vs.\ B2 as win rate and means
only (2/120 wins, ${+}1.7$\,pp; no Holm claim).
E6 (+7.7\,pp vs.\ unstructured test-only under full~M; E14 scopes
traces) remains the adoption mechanism signal.
Per-task win rate (2/120 fixed-oracle E1; E12 mean-across-seeds: 0 wins /
1 loss / 119 ties) and overlap stratification (E3; fixed-oracle E10:
M ${+}1.0$\,pp) supplement the headline means.
These align with the per-task distributions
in Figures~\ref{fig:conf-distribution} and~\ref{fig:conf-cdf}.

%=============================================================================
\section{Case Study: The Tax Bracket That Fooled the Unit Tests}
\label{sec:results:case-study}
%=============================================================================

Aggregate tables answer \emph{how often}; this section answers \emph{what it
felt like} when structured feedback actually moved a candidate.
We walk one confrontation on the running example \code{compute\_tax}
(Listings~\ref{lst:tax-fsf}--\ref{lst:tax-bug} in Chapter~\ref{ch:intro}):
foreign-high surcharge must preempt the domestic mid bracket whenever both
guards hold.
The numbers below are representative of the E6 rescue regime
(Tables~\ref{tab:e6-win-profile},~\ref{tab:e6-win-loss})---collapsed
test-only feedback, then a typed IR that names the violated scenario.

\subsection*{Act I --- The Plausible Bug}

The first LLM draft looked almost careful.
It rejected negative incomes, handled a mid bracket, and even mentioned the
foreign-high surcharge---but in the wrong order:

\inputpython{listings/compute_tax_bug.py}{First draft: mid bracket evaluated
  before foreign-high surcharge.}{lst:case-bug}

On everyday domestic incomes the function looks fine.
A hand-written unit suite that samples
$(\mathit{income},\mathit{is\_foreign})\in\{(30000,0),(60000,0),(90000,0)\}$
returns \texttt{Low}/\texttt{DomesticMid} as expected and never draws a
foreign earner above $120{,}000$.
Pass rate on that suite: $3/3$.
The bug is not ``the model cannot code''; it is that the model coded the
\emph{natural} order of if-statements instead of the \emph{contracted}
first-match order.

\subsection*{Act II --- Why Ordinary Tests Let It Slip}

Figure~\ref{fig:defect-vs-test} is exactly this geometry.
Random or developer tests concentrate on $\mathit{is\_foreign}{=}0$ and mid
incomes---the blue cloud on the left.
The fatal overlap
$\mathit{income}{\ge}120000 \wedge \mathit{is\_foreign}{=}1$ is a thin
orange slice; without an SMT witness for $\Phi_2$, the suite never opens
that door.
Unstructured test-feedback repair, if it fires at all, typically says only
``assertion failed'' or dumps raw I/O.
That is the E6 \texttt{test\_only} regime: the model is told \emph{that}
something failed, not \emph{which scenario under first-match} failed.
In several E6 wins the test-only Conf.\ collapses to~$0$ while the typed IR
still localises the residual region
(Table~\ref{tab:e6-win-profile}).

\subsection*{Act III --- Semantic Feedback IR Names the Wound}

HSP-Agile does not argue with the blue cloud.
It asks Z3 for a model of $\Phi_2$
(Equation~\eqref{eq:phi-intro}) and gets $(150000,1)$.
Oracle: \texttt{ForeignHigh}. Candidate: \texttt{DomesticMid}.
The failure is packaged as a typed record before any English is rendered
(Listing~\ref{lst:case-ir}).

\begin{lstlisting}[language={},basicstyle=\ttfamily\footnotesize,frame=single,
  caption={Semantic Feedback IR for the foreign-high miss (abbreviated).},
  label={lst:case-ir}]
{
  "violation_type": "ordering",
  "scenario_index": 2,
  "constraint_text": "income >= 120000 && is_foreign == 1",
  "inputs": {"income": 150000, "is_foreign": 1},
  "expected": {"bracket": "ForeignHigh"},
  "observed": {"bracket": "DomesticMid"},
  "reason": "Scenario S2 must preempt S3 under first-match.",
  "suggested_fix": "Evaluate foreign-high before domestic mid."
}
\end{lstlisting}

The repair prompt then says, in substance:
\begin{quote}
\textit{For inputs} $(\mathit{income}{=}150000,\;\mathit{is\_foreign}{=}1)$,
\textit{your code returned} \texttt{DomesticMid}, \textit{but the
specification requires} \texttt{ForeignHigh}.
\textit{This input satisfies Scenario~$S_2$}
(\texttt{income >= 120000 \&\& is\_foreign == 1}),
\textit{which must be matched before Scenario~$S_3$}.
\textit{Suggested fix: evaluate the foreign-high surcharge before the
domestic mid bracket.}
\end{quote}
That paragraph is the difference between E6 variants A/B and C: not more
tokens for their own sake, but scenario identity, expected value, and a
precedence hint the model can act on.

\subsection*{Act IV --- The Rewrite}

After one repair call ($T_\text{repair}{=}0$), the candidate restores
contract order:

\inputpython{listings/compute_tax_fixed.py}{Repaired candidate: foreign-high
  before domestic mid.}{lst:case-fixed}

Every SpecIR witness now matches, including $(150000,1)$; the structural
screen is clear of RF07-style unconditional success; $\mathrm{Accept}{=}1$.
Mean Conf.\ on the task jumps from a rescue-from-zero regime under
test-only feedback to full witness agreement under typed IR---the same
qualitative arc as the 13/14 E6 C$-$A wins that began with
\texttt{test\_only} Conf${=}0$.

\subsection*{What the Story Is Not}

This case is an \emph{ordering} confrontation that motivates witnesses and
the Accept gate.
The dominant E6 Conf.\ \emph{gain} on BENCH-120 is still catch-all
\texttt{others}/arithmetic rescue
(Table~\ref{tab:defect-evidence-map}); we chose \code{compute\_tax} here
because the running example makes the missed overlap legible without a
synthetic HardCase id.
The moral is the same either way: unstructured pass/fail leaves the model
guessing; Semantic Feedback IR points at the wound and asks for a rewrite.

\bigskip


\section{Historical Failure Analysis (E9 Archive)}
\label{sec:app:e9}
% label on Chapter~8 pointer: sec:disc:failure-analysis
%=============================================================================

\paragraph{Archive status.}
E9 was collected under the \emph{pre-fix} others-witness oracle, where
\textbf{114 of 120 tasks (95.0\%)} failed conjunctive acceptance after
$K{=}3$.
Under fixed-oracle strengthened~M (\texttt{run\_e1\_m\_win\_v2}), Strict for~M
is 100.0\%---so the E9 pie describes \textbf{historical residual mass only},
not the current primary Strict rate.
We retain it as a qualitative catalogue of failure \emph{modes} (ordering,
boundary, arithmetic) that still guide future witness/IR design; we do
\emph{not} cite E9 percentages as current effectiveness evidence.
(Re-running E9 under the fixed oracle is future work: near-ceiling Strict
leaves too few residuals for a powered taxonomy.)

Despite strong mean conformance on several tiers, the historical E9 residual
taxonomy assigned a primary failure category to \textbf{90} of 114
non-accepted tasks under the pre-fix oracle.
The pie in Figure~\ref{fig:failure-taxonomy} maps that archive residual mass.

\begin{figure}[t]
  \centering
  \thesisfig[max width=0.72\linewidth]{figures/failure_taxonomy_pie.pdf}
  \caption{\textbf{Claim (historical / pre-fix oracle):} residual failures
    after $K=3$ on \texttt{run\_e1\_canonical\_v1} are dominated by ordering
    and boundary errors, not random noise.
    Under fixed-oracle E1, M Strict is 100.0\%---this pie is not the current
    primary residual mass.
    \textbf{How to read:} slices are shares of the $n{=}90$ E9-labelled
    non-accepted tasks under mode~M (of 114 total non-accepted).}
  \label{fig:failure-taxonomy}
\end{figure}

Of the 90 E9-labelled non-accepted tasks after $K{=}3$: OrderingError 31.1\% (LLM re-asserts
wrong order), BoundaryError 26.7\% (off-by-one), ArithmeticError 17.8\%,
StateError 13.3\%, Hallucination 5.6\% (often RF07), MissingConstraint /
Other the rest.
Implication: richer ordering/arithmetic feedback and denser boundary
witnesses target the majority; multi-step witnesses target StateError.

%-----------------------------------------------------------------------------
\subsection{Named Failure Patterns}
\label{sec:disc:failure-patterns}
%-----------------------------------------------------------------------------

Five recurring failure patterns account for the bulk of the 90 residual errors.

\paragraph{Arithmetic reasoning failure.}
The LLM produces a formula that is structurally correct (correct branches,
correct variable references) but arithmetically wrong---for example, computing a
weighted sum with the wrong denominator or using integer division where float
division is required.
HSP-Agile misses these because Z3 witnesses are generated from guard predicates,
not from output expressions; the formal checker therefore cannot witness an
expression-level arithmetic fault unless it happens to coincide with a guard
boundary.

\paragraph{Long-range dependency.}
Some FSF specifications require the output of scenario~$i$ to depend on the
\emph{cumulative} effect of scenarios $1, \ldots, i-1$ (e.g.\ a running total or
a sequence-sensitive state machine transition).
The repair loop issues one witness per scenario in isolation; it cannot construct
a multi-step witness that would expose the accumulated-state fault.
This structural limitation of the E9 formal checker (single-scenario witnesses)
is the primary driver of \emph{OrderingError} and \emph{StateError} failures.

\paragraph{State update confusion.}
When two or more FSF scenarios modify the same output variable, the LLM
occasionally generates code in which later scenarios overwrite rather than
accumulate the earlier value.
The counterexample-guided repair prompt identifies the violated scenario but does
not explicitly communicate the accumulation invariant, so the LLM corrects the
reported case while re-introducing the error in an adjacent scenario.
Extending the Semantic Feedback IR with an explicit \texttt{accumulation\_mode}
field would address this pattern.

\paragraph{Unsatisfiable guard (Unsat guard).}
A small number of FSF specifications contain scenario guards that Z3 proves
mutually satisfiable but that the LLM interprets as mutually exclusive---producing
dead branches or guard conditions that can never be simultaneously true.
Because the implementation is syntactically valid and passes individual-scenario
witnesses, the formal checker does not flag the dead branch; only a path-coverage
oracle (not currently part of the pipeline) would detect it.

\paragraph{Hallucinated API call.}
The LLM occasionally calls a library function that does not exist in the
permitted API surface defined by the FSF signature---for example, calling
\texttt{math.log\_base(x, b)} instead of \texttt{math.log(x) / math.log(b)}.
Pattern guard RF07 catches this post-generation, but the repair loop does not
receive a structured signal about the illegal call; it receives only a generic
runtime error counterexample, leading to repeated hallucinations across repair
attempts.
A dedicated RF07-triggered repair template that injects the legal API surface into
the repair prompt would break this cycle.

\subsection{Failure Trace: Scenario Ordering}
\label{subsec:failure-trace-ordering}

The failed \code{classify\_signal} candidate evaluates Scenario~2 before
Scenario~1 (Listing~\ref{lst:llm-incorrect}, Section~\ref{sec:intro:problem}).
Witness $\mathbf{x}^*=(\texttt{level}{=}{-3},\,\texttt{threshold}{=}{-10})$
activates $S_1$ under first-match semantics; the checker emits:

\begin{lstlisting}[language={},caption={Semantic Feedback IR (iteration 1).},
  label=lst:classify-signal-ir]
{
  "violation_type": "ordering",
  "scenario_index": 1,
  "constraint_text": "level < 0",
  "inputs": {"level": -3, "threshold": -10},
  "expected": {"status": "Error"},
  "observed": {"status": "Critical"},
  "reason": "Scenario 1 violated under first-match guard precedence.",
  "priority": 1,
  "suggested_fix": "reorder guard evaluation to match specification precedence"
}
\end{lstlisting}

\texttt{FeedbackRenderer} projects this JSON record into the repair prompt;
at $k=2$ the candidate fixes the negative-threshold branch but re-introduces
an ordering error on a different overlap region---illustrating why E9
\emph{OrderingError} remains the largest residual category despite structured IR.

\subsection{Failure Trace: Boundary Reasoning}
\label{subsec:failure-trace}

Task \texttt{HardCase036} specifies five scenarios with overlapping boundary
guards.  At repair iteration $k=2$, witness $\mathbf{x}^*=(a=2,b=6,c=3)$
activates Scenario~1 under first-match semantics.  The Semantic Feedback IR is:

\begin{lstlisting}[language={},caption={Semantic Feedback IR (iteration 2).},
  label=lst:failure-ir]
{
  "violation_type": "ordering",
  "scenario_index": 1,
  "constraint_text": "a le 3 && b gt 5",
  "inputs": {"a": 2, "b": 6, "c": 3},
  "expected": {"risk": 2, "action": 1},
  "observed": {"risk": 1, "action": 3},
  "reason": "guard evaluated out of order",
  "priority": 1,
  "suggested_fix": "reorder"
}
\end{lstlisting}

Despite structured IR at $k=2$, repair at $k=3$ introduces a new ordering
violation, illustrating \emph{long-range dependency} beyond single-step repair.

\paragraph{Guard-graph pattern matching (future).}
Representing SpecIR cases as a guard dependency graph enables reachability-based
detection of shadow guards and unreachable paths beyond the current PoC
smell catalogue---or by swapping in a stronger pluggable $\mathrm{Screen}$.


%=============================================================================
